\chapter*{Введение}
\addcontentsline{toc}{chapter}{Введение}

В современных организациях значительная часть прикладных знаний хранится в документах: регламентах, инструкциях, технических описаниях, коммерческих материалах, таблицах, ответах на типовые вопросы и методических рекомендациях. По мере накопления таких материалов возрастает сложность поиска нужной информации. Пользователь может не знать точного названия файла, формулировать вопрос естественным языком, использовать обобщенные термины или ожидать не ссылку на документ, а краткий ответ, пригодный для практического применения.

Корпоративная база знаний RobyMAX в рамках рассматриваемого проекта представлена русскоязычными бизнесово-техническими документами, связанными с продуктами, техническими требованиями, пресейлом, внедрением, сервисом и коммерческими условиями. Такие документы неоднородны по формату и структуре: часть материалов представлена в DOCX, часть — в PDF, часть — в XLSX. Ответ на вопрос может находиться не в названии документа, а внутри раздела, таблицы или фрагмента текста. Поэтому задача поиска по такой базе знаний не сводится к просмотру списка файлов.

Ручной поиск по папкам и документам требует времени и знания структуры архива. Полнотекстовый поиск снижает трудоемкость, но в основном ориентирован на совпадение слов и выражений. В корпоративной среде это ограничение существенно: один и тот же смысл может быть выражен разными формулировками, а пользовательский вопрос часто задается в прикладной форме. Например, вопрос о требованиях к инфраструктуре может соответствовать фрагментам про сеть Wi-Fi, маршруты движения, доступность помещения или условия пилотных испытаний. В такой ситуации важно находить не только документы с совпадающими словами, но и фрагменты, близкие по смыслу.

Развитие больших языковых моделей создало предпосылки для построения вопросно-ответных ассистентов, способных формировать связные ответы на естественном языке~\cite{Vaswani2017Attention}. Однако языковая модель сама по себе опирается преимущественно на параметрические знания, полученные во время обучения. Для внутренних документов организации это недостаточно: модель может не знать актуального содержимого базы знаний, ошибаться в деталях или формировать правдоподобные, но неподтвержденные утверждения. Поэтому для прикладной корпоративной системы требуется связать генерацию ответа с проверяемыми источниками.

Подход Retrieval-Augmented Generation объединяет поиск релевантного контекста и генерацию ответа по найденным фрагментам~\cite{Lewis2020RAG}. Документы предварительно обрабатываются, разбиваются на фрагменты, кодируются embedding-моделью и индексируются в векторной базе данных. При поступлении вопроса система строит embedding вопроса, находит близкие чанки, формирует контекст и передает его языковой модели. Такой подход позволяет обновлять базу знаний без дообучения LLM, возвращать пользователю источники и уменьшать риск неподтвержденных ответов.

Актуальность работы определяется необходимостью построения воспроизводимого прототипа вопросно-ответного ассистента для русскоязычной корпоративной базы знаний. Новизна работы в рамках учебно-исследовательского этапа состоит не в разработке нового алгоритма машинного обучения, а в инженерном исследовании применимости базового RAG-подхода к неоднородному корпусу документов RobyMAX и в подготовке минимальной архитектуры, которую можно последовательно развивать на следующих этапах.

Целью работы является исследование и разработка базового прототипа RAG-ассистента, обеспечивающего поиск релевантных фрагментов в корпусе русскоязычных бизнесово-технических документов RobyMAX и генерацию ответа на пользовательский вопрос с использованием найденного контекста.

Для достижения цели были поставлены следующие задачи:
\begin{compactenum}
  \item провести анализ проблемы поиска информации в корпоративной базе знаний;
  \item рассмотреть существующие подходы: ручной поиск, полнотекстовый поиск, embedding retrieval и RAG;
  \item обосновать выбор базового RAG-подхода для первой MVP-версии;
  \item подготовить корпус документов для обработки;
  \item реализовать парсинг документов и приведение текста к единому структурированному виду;
  \item выполнить разбиение документов на смысловые фрагменты;
  \item построить векторные представления чанков;
  \item организовать хранение и поиск чанков в Qdrant;
  \item реализовать базовый сценарий генерации ответа по найденному контексту;
  \item выполнить первичную проверку работоспособности прототипа.
\end{compactenum}

Объектом работы является процесс поиска информации в русскоязычной корпоративной базе знаний. Предметом работы являются методы построения базового RAG-пайплайна, включающего подготовку документов, векторный поиск и генерацию ответа по найденному контексту.

Пояснительная записка состоит из четырех глав. В первой главе анализируются предметная область RobyMAX, проблема поиска информации в документах и подходы к построению вопросно-ответных систем. Во второй главе рассматриваются теоретические основы LLM, embeddings, векторных баз данных и Retrieval-Augmented Generation. В третьей главе описывается проектирование базового RAG-ассистента, включая требования, архитектуру обработки документов, поиска и генерации ответа. В четвертой главе приводится описание реализации MVP, сценария запуска, используемых инструментов, первичной проверки и ограничений полученного прототипа.

%%% Local Variables:
%%% TeX-engine: xetex
%%% End:
